% !TEX program = xelatex
% Biên dịch: xelatex finsight_ke_hoach.tex  (chạy 2 lần để cập nhật mục lục)
\documentclass[11pt,a4paper]{article}

\usepackage{fontspec}        % XeLaTeX: hỗ trợ UTF-8 tiếng Việt sẵn (font Latin Modern)
\usepackage[a4paper,margin=2.3cm]{geometry}
\usepackage{enumitem}
\usepackage{booktabs}
\usepackage{longtable}
\usepackage{array}
\usepackage{xcolor}
\usepackage{tcolorbox}
\usepackage{fancyhdr}
\usepackage{titlesec}
\usepackage[hidelinks]{hyperref}

\definecolor{fsblue}{RGB}{20,80,140}
\definecolor{fsgray}{RGB}{90,90,90}
\definecolor{fslight}{RGB}{238,243,250}

\hypersetup{colorlinks=true, linkcolor=fsblue, urlcolor=fsblue}

\titleformat{\section}{\Large\bfseries\color{fsblue}}{\thesection}{0.6em}{}
\titleformat{\subsection}{\large\bfseries\color{fsblue!85!black}}{\thesubsection}{0.6em}{}

\setlist[itemize]{leftmargin=1.4em, itemsep=2pt, topsep=2pt}
\setlist[enumerate]{leftmargin=1.6em, itemsep=2pt, topsep=2pt}

\newcommand{\pha}[1]{\subsection{#1}}
\newtcolorbox{ghichu}{colback=fslight, colframe=fsblue, boxrule=0.6pt, arc=2pt,
  left=6pt, right=6pt, top=4pt, bottom=4pt}

\pagestyle{fancy}
\fancyhf{}
\lhead{\small\color{fsgray} FinSight}
\rhead{\small\color{fsgray} Kế hoạch dự án}
\cfoot{\small\thepage}
\renewcommand{\headrulewidth}{0.3pt}

\title{\vspace{-1.2cm}\textbf{\color{fsblue}FinSight}\\[4pt]
\large Hệ thống Phân tích Báo cáo Tài chính Đa năm bằng AI Agent\\[6pt]
\normalsize\color{fsgray} Bản đặc tả kế hoạch dự án --- Tiếng Việt}
\author{}
\date{}

\begin{document}
\maketitle
\thispagestyle{fancy}
\vspace{-1.0cm}

\begin{ghichu}
\textbf{Một câu mô tả.} FinSight nhận Báo cáo tài chính (BCTC) của một doanh nghiệp qua nhiều năm (PDF hoặc bản scan tiếng Việt), \textbf{OCR giữ nguyên cấu trúc} $\rightarrow$ \textbf{trích xuất toàn bộ thông tin} $\rightarrow$ dựng \textbf{đồ thị tri thức nối các năm (GraphRAG)} $\rightarrow$ dùng \textbf{AgentRAG} (agent, tool-calling, memory) để tự phân tích, đối chiếu chéo và trả lời các nhiệm vụ cụ thể, có trích dẫn. Con người có thẩm quyền duyệt; phản hồi của họ quay lại tinh chỉnh mô hình.
\end{ghichu}

\tableofcontents
\vspace{0.4cm}

% =====================================================================
\section{Tổng quan dự án}

\subsection{Mục tiêu}
Xây dựng một hệ thống \emph{tự động phân tích} BCTC tiếng Việt: thay vì một con chatbot hỏi--đáp thụ động, hệ thống \textbf{chủ động sinh ra bản phân tích} ngay khi nạp dữ liệu, sau đó cho phép người có thẩm quyền đào sâu bằng hỏi--đáp và phê duyệt kết quả.

\subsection{Luồng hoạt động cốt lõi}
\begin{verbatim}
   BCTC nhieu nam (PDF/scan)
            |
   [1] OCR giu cau truc  -> document model co cau truc
            |
   [2] Trich xuat + chuan hoa  -> so lieu tai chinh chuan (TT200)
            |
   [3] GraphRAG  -> do thi tri thuc noi cac nam voi nhau
            |
   [4] PHAN TICH TU DONG (push) -> ty so, Altman Z, Beneish M, reconcile
            |
   [5] AgentRAG (agent + tool + memory) -> nhiem vu cu the, verify, cite
            |
   [6] Nguoi co tham quyen DUYET / sua  (HITL)
            |
   [7] Feedback -> finetune  (model gioi len)
\end{verbatim}

\begin{ghichu}
\textbf{Phân biệt quan trọng.} \emph{Trích xuất} và \emph{phân tích} KHÔNG phải RAG: trích xuất là OCR + mô hình bảng; phân tích là công thức tài chính + ML trên số đã trích. \textbf{RAG/GraphRAG} chỉ phục vụ bước \emph{nối dữ liệu xuyên năm}, \emph{kiểm chứng với quy định} và \emph{hỏi--đáp}.
\end{ghichu}

\subsection{Giá trị cho doanh nghiệp (ngân hàng / fintech)}
\begin{itemize}
  \item Rút ngắn thời gian thẩm định 1 bộ BCTC từ vài giờ xuống vài phút, chuẩn hóa.
  \item Giám sát liên tục theo thời gian, cảnh báo sớm khi tình hình xấu đi.
  \item Phát hiện dấu hiệu xào nấu sổ sách (Beneish M-Score, lệch lợi nhuận -- dòng tiền) trước khi cho vay.
  \item Mọi kết luận đều có \textbf{trích dẫn nguồn} $\rightarrow$ đáp ứng yêu cầu kiểm toán / tuân thủ.
\end{itemize}

\subsection{Phạm vi giai đoạn đầu (MVP)}
\begin{itemize}
  \item \textbf{Một công ty} niêm yết, BCTC \textbf{năm}, \textbf{5 năm} gần nhất, bản hợp nhất.
  \item Gồm: Bảng cân đối kế toán, Báo cáo KQKD, Báo cáo lưu chuyển tiền tệ, Thuyết minh.
  \item Toàn bộ \textbf{tiếng Việt}.
  \item BCTC quý sẽ bổ sung ở pha sau (mở khóa đối chiếu quý--năm và cảnh báo sớm theo tốc độ).
\end{itemize}

% =====================================================================
\section{Kiến trúc tổng thể}

\begin{verbatim}
            +---------------- INGEST SONG SONG (1 lan, bat dong bo) ----------------+
  PDF/scan  |  [P1] OCR ->  [P2] Trich xuat -> [P2] Chuan hoa (TT200)              |
  nhieu file|  (Paddle/      (TableTransformer/   (canonical chart of accounts)    |
            |   VietOCR)      Donut + LoRA)                                        |
            +----------------------------------|----------------------------------+
                                               v
                               [P3] DO THI TRI THUC (GraphRAG)
                               node: chi tieu/nam/gia tri, thuyet minh
                               edge: thuoc nam, tham chieu, year-over-year
                                               |
                  +----------------------------+----------------------------+
                  v                            v                            v
          [P4] PHAN TICH              [P5] AGENT (LangGraph)        [RAG quy dinh]
          ty so, Altman Z,            tool-calling, memory,         vbpl/TT200
          Beneish M, reconcile        nghi ngo -> verify -> cite    + web search
                  |                            |                            |
                  +----------------------------+----------------------------+
                                               v
                            [P4/P5] BAN PHAN TICH + trich dan
                                               |
              [P6] HITL duyet/sua --> corrections --> [P7] finetune LoRA
                                               |
                            [P9] Session / Memory store (moi ho so)
\end{verbatim}

% =====================================================================
\section{Tech stack}

\renewcommand{\arraystretch}{1.25}
\begin{longtable}{>{\raggedright\arraybackslash}p{3.6cm} >{\raggedright\arraybackslash}p{6.4cm} >{\raggedright\arraybackslash}p{5.0cm}}
\toprule
\textbf{Tầng} & \textbf{Công nghệ đề xuất} & \textbf{Lý do} \\
\midrule
\endhead
API / điều phối & FastAPI + Celery + Redis (hoặc Prefect/Temporal) & Bất đồng bộ, xử lý song song nhiều file \\
Đọc PDF & PyMuPDF; phát hiện text-layer vs scan & Tách đường nhanh: bỏ qua OCR khi có text-layer \\
OCR tiếng Việt & VietOCR / PaddleOCR (PP-Structure) & OCR có dấu + phân tích bố cục + bảng \\
Trích xuất bảng (DL) & Table Transformer (TATR) hoặc Donut/Qwen2-VL, finetune LoRA & Bảng đa tầng, ô gộp --- phần học sâu lõi \\
Chuẩn hóa & Mapping $\rightarrow$ mã số chỉ tiêu TT200 & Để đối chiếu \& tính tỷ số nhất quán \\
Đồ thị tri thức & Neo4j (hoặc networkx + Postgres) + GraphRAG & ``Nối các năm'' = duyệt đồ thị đa bước \\
RAG & Qdrant/pgvector + embedding bge-m3 + reranker & Truy hồi quy định, ngữ cảnh xuyên file \\
Agent & LangGraph & Luồng tường minh, có thể kiểm toán \\
LLM & API (Claude) cho điều phối + SLM finetune (distill) cho việc rẻ & Cân đối chi phí/tốc độ \\
Bất định & MAPIE / conformal prediction & Cam kết độ phủ + biết ``không chắc'' \\
Lưu trữ & Postgres(+pgvector), MinIO/S3 (file), Redis (cache) & Chuẩn production \\
Serving & vLLM cho SLM & Suy luận theo lô, nhanh \\
Quan sát & Langfuse (trace LLM) + Prometheus/Grafana & Đo độ trễ/chi phí, gỡ lỗi agent \\
HITL UI & Streamlit $\rightarrow$ Next.js & Sửa trường, duyệt cờ đỏ \\
\bottomrule
\end{longtable}

% =====================================================================
\section{Đặc tả chi tiết từng pha}

Mỗi pha mô tả theo khung: \emph{Mục tiêu --- Đầu vào --- Các bước/Thành phần --- Công nghệ --- Đầu ra --- Tiêu chí hoàn thành (DoD) --- Rủi ro}.

% --------------------------------------------------
\pha{Pha 0 --- Khởi tạo \& Thu thập dữ liệu}
\begin{itemize}
  \item \textbf{Mục tiêu:} dựng khung repo, hạ tầng dev, và bộ dữ liệu khởi đầu.
  \item \textbf{Các bước:}
  \begin{enumerate}
    \item Tải 5 BCTC năm (gồm cả bản text-layer và bản scan) của 1 công ty niêm yết từ cafef.vn / vietstock.vn / hsx.vn.
    \item Tải văn bản chuẩn: Thông tư 200/2014/TT-BTC (chế độ kế toán), TT96/2020 (công bố thông tin) từ vbpl.vn.
    \item Khởi tạo repo, môi trường (Python 3.11, Poetry/uv), \texttt{docker-compose} (Postgres, Redis, MinIO, Qdrant).
    \item Thiết lập CI cơ bản (lint, test, pre-commit).
  \end{enumerate}
  \item \textbf{Đầu ra:} repo chạy được, dữ liệu thô trong \texttt{data/raw/}.
  \item \textbf{DoD:} \texttt{docker-compose up} dựng đủ dịch vụ; 5 PDF tải về; sơ đồ schema DB phác thảo.
  \item \textbf{Rủi ro:} bản scan chất lượng kém $\rightarrow$ chuẩn bị thêm 1--2 công ty dự phòng.
\end{itemize}

% --------------------------------------------------
\pha{Pha 1 --- OCR giữ nguyên cấu trúc}
\begin{itemize}
  \item \textbf{Mục tiêu:} biến PDF/scan thành \emph{document model có cấu trúc}, không mất bố cục bảng.
  \item \textbf{Đầu vào:} PDF (text-layer hoặc ảnh scan).
  \item \textbf{Các bước/Thành phần:}
  \begin{enumerate}
    \item Phát hiện loại trang: có text-layer $\rightarrow$ đọc trực tiếp; là ảnh $\rightarrow$ đưa qua OCR.
    \item Phân tích bố cục (layout analysis): nhận diện vùng bảng, đoạn văn, tiêu đề, header/footer, số trang.
    \item OCR tiếng Việt cho vùng text + ô bảng (giữ tọa độ bounding box).
    \item Tái dựng cấu trúc: bảng $\rightarrow$ lưới ô (hàng/cột, ô gộp); văn bản $\rightarrow$ thứ tự đọc đúng.
  \end{enumerate}
  \item \textbf{Công nghệ:} PyMuPDF; PaddleOCR PP-Structure (layout + table) ; VietOCR (có thể finetune cho phông chữ BCTC).
  \item \textbf{Đầu ra:} JSON document model: \texttt{\{pages:[\{blocks, tables:[\{cells, bbox\}], reading\_order\}]\}}.
  \item \textbf{Đánh giá:} CER/WER cho text; \textbf{TEDS} (Tree-Edit-Distance Similarity) cho cấu trúc bảng.
  \item \textbf{DoD:} 1 BCTC scan $\rightarrow$ JSON cấu trúc đầy đủ, bảng cân đối tái dựng đúng $\geq 95\%$ ô.
  \item \textbf{Rủi ro:} ô gộp/đa tầng header $\rightarrow$ cần finetune + hậu xử lý theo luật.
\end{itemize}

% --------------------------------------------------
\pha{Pha 2 --- Trích xuất toàn bộ thông tin \& Chuẩn hóa}
\begin{itemize}
  \item \textbf{Mục tiêu:} từ document model $\rightarrow$ số liệu tài chính \emph{có cấu trúc, chuẩn hóa} (BCTC rất dài).
  \item \textbf{Các bước:}
  \begin{enumerate}
    \item Nhận dạng cấu trúc bảng (table structure recognition) cho 4 báo cáo chính + các bảng trong thuyết minh.
    \item Trích metadata: kỳ báo cáo, đơn vị tiền (nghìn/triệu đồng), tên công ty, mã số thuế, trạng thái kiểm toán.
    \item Parse số: dấu phân tách nghìn kiểu VN, số âm trong ngoặc \texttt{(1.234)}, ô trống.
    \item \textbf{Chuẩn hóa:} ánh xạ mỗi dòng $\rightarrow$ \textbf{mã số chỉ tiêu TT200} (ví dụ ``Tiền và tương đương tiền'' $\rightarrow$ mã 110). Đây là cầu nối cho phân tích \& đối chiếu.
  \end{enumerate}
  \item \textbf{Công nghệ:} Table Transformer / Donut + \textbf{LoRA finetune} trên bảng BCTC tiếng Việt; rule-based normalizer.
  \item \textbf{Đầu ra:} bảng số liệu chuẩn hóa theo năm trong Postgres (schema canonical).
  \item \textbf{Schema 1 dòng (ví dụ):}
  \begin{verbatim}
  { "chi_tieu": "Tien va tuong duong tien", "ma_so": "110",
    "thuyet_minh": "V.01", "nam": 2024,
    "gia_tri": 12345678901, "don_vi": "VND" }
  \end{verbatim}
  \item \textbf{Đánh giá:} độ chính xác theo ô (đặc biệt ô số), độ chính xác ánh xạ dòng $\rightarrow$ mã số.
  \item \textbf{DoD:} 5 năm $\rightarrow$ bảng chuẩn hóa đầy đủ; golden set $\sim$15--20 bảng để đo.
  \item \textbf{Rủi ro:} tên chỉ tiêu khác nhau giữa các năm/công ty $\rightarrow$ cần từ điển đồng nghĩa + fuzzy match.
\end{itemize}

% --------------------------------------------------
\pha{Pha 3 --- Đồ thị tri thức đa năm (GraphRAG)}
\begin{itemize}
  \item \textbf{Mục tiêu:} ``nối các năm với nhau'' để truy vấn đa bước, đối chiếu xuyên file.
  \item \textbf{Thiết kế đồ thị:}
  \begin{itemize}
    \item \textbf{Node:} chỉ tiêu (mã số), giá trị theo năm, công ty, mục thuyết minh, kỳ báo cáo.
    \item \textbf{Edge:} \emph{thuộc năm}, \emph{thuộc báo cáo}, \emph{tham chiếu thuyết minh}, \emph{quan hệ năm-qua-năm}, \emph{quan hệ tính toán} (ví dụ Lợi nhuận gộp = Doanh thu $-$ Giá vốn).
  \end{itemize}
  \item \textbf{GraphRAG indexing:} trích entity + quan hệ, sinh tóm tắt cụm (community summary) để trả lời câu hỏi ``toàn cục''.
  \item \textbf{Công nghệ:} Neo4j (hoặc networkx + Postgres); GraphRAG (Microsoft hoặc tự xây); embedding bge-m3 cho truy hồi.
  \item \textbf{Đầu ra:} đồ thị truy vấn được; API \texttt{query\_graph()} trả dữ liệu xuyên năm.
  \item \textbf{DoD:} truy vấn ``doanh thu 2020--2024'' và ``đường tính lợi nhuận gộp năm 2023'' chạy đúng.
  \item \textbf{Rủi ro:} bùng nổ node $\rightarrow$ chỉ mô hình hóa chỉ tiêu cốt lõi + thuyết minh quan trọng.
\end{itemize}

% --------------------------------------------------
\pha{Pha 4 --- Phân tích tự động \& Đối chiếu (push)}
\begin{itemize}
  \item \textbf{Mục tiêu:} ngay khi có dữ liệu, \emph{chủ động} sinh bản phân tích, không chờ hỏi.
  \item \textbf{Tính toán:}
  \begin{itemize}
    \item \textbf{Tỷ số \& xu hướng:} thanh khoản (current, quick), đòn bẩy (D/E, Nợ/EBITDA), sinh lời (ROA, ROE, biên LN), khả năng trả lãi (interest coverage), vòng quay vốn (DSO/DIO/DPO).
    \item \textbf{Dự báo phá sản:} Altman Z'' (cho doanh nghiệp dùng giá trị sổ sách).
    \item \textbf{Phát hiện xào nấu:} Beneish M-Score (8 biến, cần 2 năm liên tiếp); accruals; lệch lợi nhuận -- dòng tiền; Benford.
    \item \textbf{Đối chiếu (reconciliation):} thuyết minh vs số trên báo cáo; nhất quán năm-qua-năm; (sau: lũy kế quý vs năm).
  \end{itemize}
  \item \textbf{Bất định:} dùng conformal prediction để \emph{abstain} (``không đủ cơ sở'') thay vì đoán bừa.
  \item \textbf{Đầu ra:} \textbf{Bản phân tích / Thẻ rủi ro} có cờ đỏ + trích dẫn vị trí số liệu.
  \item \textbf{DoD:} sinh thẻ phân tích tự động cho 1 công ty 5 năm, mỗi cờ đỏ truy được về dòng nguồn.
  \item \textbf{Rủi ro:} công thức sai do đơn vị/dấu $\rightarrow$ unit test cho từng chỉ số.
\end{itemize}

% --------------------------------------------------
\pha{Pha 5 --- AgentRAG: Agent, Tool-calling, Memory}
\begin{itemize}
  \item \textbf{Mục tiêu:} cho phép giao \emph{nhiệm vụ cụ thể} bằng ngôn ngữ tự nhiên trên nền dữ liệu đã phân tích.
  \item \textbf{Các loại nhiệm vụ (theo đúng mô tả):}
  \begin{itemize}
    \item ``Trích xuất toàn bộ thông tin BCTC năm 2023.''
    \item ``Phần nào khớp / phần nào lệch qua các năm?'' (nhiệm vụ đối chiếu).
    \item ``So sánh khoản mục \emph{hàng tồn kho} giữa 2020--2024.''
    \item ``Tính và giải thích biến động biên lợi nhuận gộp 5 năm.''
  \end{itemize}
  \item \textbf{Tools (function calling):} \texttt{query\_graph}, \texttt{get\_lineitem}, \texttt{compute\_ratio}, \texttt{reconcile}, \texttt{search\_regulation} (RAG vbpl/TT200), \texttt{web\_search}.
  \item \textbf{Memory:}
  \begin{itemize}
    \item Ngắn hạn: ngữ cảnh hội thoại hiện tại.
    \item Dài hạn: phát hiện đã được duyệt, lưu theo \emph{session = một hồ sơ doanh nghiệp}.
  \end{itemize}
  \item \textbf{Verify khi nghi ngờ:} confidence thấp / nghi vi phạm $\rightarrow$ truy hồi quy định + (nếu cần) web search nguồn chính thống $\rightarrow$ \textbf{trích dẫn}; abstain khi không đủ bằng chứng.
  \item \textbf{Công nghệ:} LangGraph (đồ thị quyết định tường minh, audit được).
  \item \textbf{Đầu ra:} câu trả lời/nhiệm vụ hoàn thành kèm trích dẫn + đường suy luận lưu lại.
  \item \textbf{DoD:} agent hoàn thành đúng 4 loại nhiệm vụ mẫu, mỗi kết quả có trích dẫn.
  \item \textbf{Rủi ro:} ảo giác số liệu $\rightarrow$ \emph{bắt buộc} agent chỉ lấy số qua tool, không tự ``nhớ'' số.
\end{itemize}

% --------------------------------------------------
\pha{Pha 6 --- Human-in-the-loop \& Vòng tinh chỉnh}
\begin{itemize}
  \item \textbf{Mục tiêu:} người có thẩm quyền duyệt bản phân tích; phản hồi thành dữ liệu huấn luyện.
  \item \textbf{Các bước:}
  \begin{enumerate}
    \item UI duyệt/sửa trường trích xuất \& cờ đỏ (Streamlit).
    \item Active learning: ưu tiên đưa case confidence thấp cho người xem.
    \item Gom correction $\rightarrow$ LoRA finetune lại extractor; \textbf{trigger theo drift} (độ phủ conformal tụt / PSI), không theo lịch cứng.
  \end{enumerate}
  \item \textbf{Đầu ra:} mỗi vòng feedback cải thiện độ chính xác đo được.
  \item \textbf{DoD:} chạy trọn 1 vòng: sửa $\rightarrow$ finetune $\rightarrow$ accuracy tăng có số liệu.
  \item \textbf{Rủi ro:} dữ liệu correction ít $\rightarrow$ gộp với dữ liệu tổng hợp (synthetic) để đủ tín hiệu.
\end{itemize}

% --------------------------------------------------
\pha{Pha 7 --- Hiệu năng \& Xử lý song song nhiều file}
\begin{itemize}
  \item \textbf{Nguyên tắc:} làm nặng \emph{một lần} lúc nạp; truy vấn chỉ duyệt thứ đã dựng sẵn.
  \item \textbf{Song song 3 mức:} theo file (mỗi file 1 task Celery), theo trang (OCR song song), theo lô GPU (gộp lời gọi mô hình).
  \item \textbf{Pipeline không rào chắn:} file nào xong xử lý tiếp file đó (wall-clock = file chậm nhất).
  \item \textbf{Cache:} hash file (idempotency), cache embedding, \textbf{semantic cache} cho câu hỏi lặp.
  \item \textbf{Serving:} vLLM cho SLM; \emph{fast-path} bằng SLM rẻ, chỉ escalate model lớn khi confidence thấp.
  \item \textbf{DoD:} xử lý $N$ file song song; báo cáo throughput (file/phút) + p95 latency truy vấn.
  \item \textbf{Rủi ro:} nghẽn GPU $\rightarrow$ hàng đợi + backpressure.
\end{itemize}

% --------------------------------------------------
\pha{Pha 8 --- Đánh giá \& Backtest}
\begin{itemize}
  \item \textbf{Metrics theo tầng:} OCR (CER/WER), bảng (TEDS), trích xuất (cell accuracy), ánh xạ mã số (accuracy), đối chiếu (tỷ lệ bắt lỗi trên bộ synthetic cấy lỗi), rủi ro (backtest).
  \item \textbf{Backtest giá trị thật:} chạy Altman Z / Beneish M trên công ty đã \emph{hủy niêm yết / scandal kế toán} $\rightarrow$ chứng minh hệ thống flag được \emph{trước} sự kiện.
  \item \textbf{DoD:} báo cáo đánh giá đầy đủ + ít nhất 1 ca backtest thành công.
\end{itemize}

% --------------------------------------------------
\pha{Pha 9 --- Quan sát, Bảo mật, Triển khai}
\begin{itemize}
  \item \textbf{Quan sát:} Langfuse (trace agent/LLM, chi phí token), Prometheus + Grafana (hệ thống).
  \item \textbf{Bảo mật \& tuân thủ:} không dùng dữ liệu PII thật; audit log mọi kết luận; phân quyền người duyệt.
  \item \textbf{Session/Memory store:} mỗi hồ sơ là một phiên có thể mở lại, kèm lịch sử phát hiện \& phê duyệt.
  \item \textbf{Triển khai:} \texttt{docker-compose} cho dev; tách service (ingest, agent, UI).
  \item \textbf{DoD:} chạy end-to-end có giám sát, mở lại được session cũ.
\end{itemize}

% =====================================================================
\section{Bổ sung những phần còn thiếu (so với mô tả ban đầu)}

Mô tả gốc tập trung OCR $\rightarrow$ trích xuất $\rightarrow$ AgentRAG. Để thành một dự án ``thật sự hữu ích'' và \emph{bảo vệ được trước người chấm khắt khe}, cần bổ sung:

\renewcommand{\arraystretch}{1.2}
\begin{longtable}{>{\raggedright\arraybackslash}p{4.6cm} >{\raggedright\arraybackslash}p{10.4cm}}
\toprule
\textbf{Hạng mục bổ sung} & \textbf{Vì sao cần} \\
\midrule
\endhead
Chuẩn hóa TT200 & Cầu nối để phân tích \& đối chiếu nhất quán; thiếu nó thì số liệu rời rạc. \\
Tầng phân tích rủi ro & Biến số thô thành \emph{tín hiệu ra quyết định} (Altman, Beneish) --- đây mới là giá trị. \\
Bộ đánh giá (eval harness) & Không đo được thì không chứng minh được ``hữu ích''. \\
HITL + vòng finetune & Biến phản hồi người dùng thành cải thiện model (flywheel). \\
Conformal / abstention & Cam kết độ tin cậy; biết ``không chắc'' thay vì bịa. \\
Drift detection & Trigger huấn luyện lại đúng lúc (MLOps thật, không phải cron). \\
Quan sát (observability) & Gỡ lỗi agent, theo dõi chi phí/độ trễ. \\
Bảo mật / PII & Dữ liệu tài chính nhạy cảm; bắt buộc với ngành. \\
Hiệu năng / song song & Yêu cầu thực tế khi xử lý nhiều file. \\
Data \& model versioning & Tái lập kết quả (DVC / MLflow). \\
Kiểm thử (testing) & Unit test cho công thức tài chính, integration test pipeline. \\
\bottomrule
\end{longtable}

% =====================================================================
\section{Cấu trúc thư mục repo đề xuất}
\begin{verbatim}
finsight/
  data/                # raw/ , processed/ , golden/ , synthetic/
  configs/             # cau hinh model, pipeline, prompt
  finsight/
    ingest/            # P1: doc PDF, phat hien scan, hang doi
    ocr/               # P1: OCR + layout (Paddle/VietOCR)
    extraction/        # P2: table model + LoRA + normalizer TT200
    graph/             # P3: dung GraphRAG, query do thi
    analysis/          # P4: ty so, Altman, Beneish, reconcile
    agent/             # P5: LangGraph, tools, memory
    rag/               # RAG quy dinh (vbpl), embedding, reranker
    hitl/              # P6: review UI + thu thap correction
    serving/           # P7: vLLM, cache, API
    eval/              # P8: metrics, backtest, golden set
    common/            # schema, db, logging, config
  ui/                  # Streamlit / Next.js
  docker/              # docker-compose, Dockerfile
  tests/
  notebooks/
  README.md
\end{verbatim}

% =====================================================================
\section{Lộ trình thời gian (gợi ý 12 tuần)}
\renewcommand{\arraystretch}{1.2}
\begin{longtable}{>{\raggedright\arraybackslash}p{1.8cm} >{\raggedright\arraybackslash}p{8.2cm} >{\raggedright\arraybackslash}p{4.6cm}}
\toprule
\textbf{Tuần} & \textbf{Công việc} & \textbf{Kết quả} \\
\midrule
\endhead
1 & Pha 0: data + repo + hạ tầng & Khung chạy \\
2--3 & Pha 1--2: OCR + trích xuất text-layer + chuẩn hóa TT200 & MVP trích xuất, golden set \\
4 & Pha 4: tầng phân tích (tỷ số + Altman + Beneish) & \textbf{Đã có giá trị thật} \\
5--6 & Pha 1: OCR scan + Pha 6: LoRA finetune từ correction & Học sâu + flywheel \\
7--8 & Pha 3: GraphRAG + Pha 4: đối chiếu xuyên năm & Liên thông các năm \\
9--10 & Pha 5: AgentRAG (tool + memory) + RAG quy định + conformal & Agent + tuân thủ \\
11--12 & Pha 7--9: song song, quan sát, backtest, viết blog/README & Demo + bằng chứng \\
\bottomrule
\end{longtable}

\begin{ghichu}
\textbf{Quy tắc thực thi.} Làm \emph{lát cắt dọc} trước (trích xuất $\rightarrow$ thẻ rủi ro chạy được ở tuần 4), rồi mới đào sâu từng tầng. \emph{Không} dựng hết hạ tầng rồi mới có cái chạy.
\end{ghichu}

% =====================================================================
\section{Tiêu chí thành công}
\begin{itemize}
  \item Trích xuất: cell accuracy $\geq 95\%$ trên golden set; ánh xạ mã số $\geq 90\%$.
  \item Đối chiếu: bắt $\geq 90\%$ lỗi cấy trên bộ synthetic có ground-truth.
  \item Phân tích: backtest Altman/Beneish flag đúng $\geq 1$ ca công ty đã vỡ/scandal.
  \item Agent: hoàn thành 4 loại nhiệm vụ mẫu, 100\% kết luận số liệu có trích dẫn (không bịa số).
  \item Hiệu năng: xử lý song song, có số liệu throughput \& p95 latency.
\end{itemize}

% =====================================================================
\section{Rủi ro \& Giảm thiểu}
\renewcommand{\arraystretch}{1.2}
\begin{longtable}{>{\raggedright\arraybackslash}p{6.2cm} >{\raggedright\arraybackslash}p{8.6cm}}
\toprule
\textbf{Rủi ro} & \textbf{Giảm thiểu} \\
\midrule
\endhead
Bản scan kém, bảng phức tạp & Finetune OCR + hậu xử lý theo luật; chọn công ty BCTC sạch trước. \\
Ảo giác số liệu của LLM & Agent chỉ lấy số qua tool; cấm ``nhớ'' số; trích dẫn bắt buộc. \\
Thiếu dữ liệu correction để finetune & Bổ sung dữ liệu tổng hợp có ground-truth điều khiển được. \\
Over-claim novelty & Pitch bằng cam kết/độ đo (coverage, accuracy), không bằng tính từ. \\
Scope phình to & Khóa MVP: 1 công ty, BCTC năm, 5 năm; tính năng khác để pha sau. \\
\bottomrule
\end{longtable}

\vspace{0.4cm}
\begin{center}
\color{fsgray}\small --- Hết bản kế hoạch FinSight ---
\end{center}

\end{document}
